\section{Polar Coordinates}

\subsection{Definition}
$(r, \theta)$ with $x = r \cos\theta$, $y = r \sin\theta$. Conversion back:
$r^2 = x^2 + y^2$, $\tan\theta = y/x$ (with quadrant chosen by signs of $x, y$).

\subsection{Slope of a polar curve $r = f(\theta)$}
Treat $\theta$ as the parameter: $x = f(\theta) \cos\theta$, $y = f(\theta) \sin\theta$.
\[ \frac{dy}{dx}
= \frac{f'(\theta) \sin\theta + f(\theta) \cos\theta}{f'(\theta) \cos\theta - f(\theta) \sin\theta}. \]

\subsection{Area in polar coordinates}
The area swept by $r = f(\theta) \ge 0$ for $\theta \in [\alpha, \beta]$ (no self-overlap) is
\[ A = \tfrac{1}{2} \int_\alpha^\beta f(\theta)^2 \dth. \]
For the region between two polar curves $r = R(\theta) \ge r = r(\theta) \ge 0$:
\[ A = \tfrac{1}{2} \int_\alpha^\beta \bigl(R(\theta)^2 - r(\theta)^2\bigr) \dth. \]

\subsection{Polar arc length}
\[ L = \int_\alpha^\beta \sqrt{f(\theta)^2 + f'(\theta)^2} \dth. \]

\subsection{Useful curves}
\[
\renewcommand{\arraystretch}{1.3}
\begin{array}{l|l}
\text{Curve} & r(\theta) \\\hline
\text{Circle, radius } a, \text{centered at origin} & r = a \\
\text{Circle through origin, radius } a/2 & r = a \cos\theta \text{ or } a \sin\theta \\
\text{Cardioid} & r = a(1 + \cos\theta) \\
\text{Lima\c con} & r = a + b \cos\theta \\
\text{Rose, } n \text{ petals if odd, } 2n \text{ if even} & r = a \cos(n\theta) \\
\text{Lemniscate} & r^2 = a^2 \cos(2\theta) \\
\text{Spiral of Archimedes} & r = a\theta
\end{array}
\]

\begin{example}[Cardioid area]
$r = 1 + \cos\theta$ on $[0, 2\pi]$:
\[ A = \tfrac{1}{2}\!\int_0^{2\pi}\! (1 + \cos\theta)^2 \dth
= \tfrac{1}{2}\!\int_0^{2\pi}\! \!\Bigl(1 + 2\cos\theta + \tfrac{1 + \cos 2\theta}{2}\Bigr) \dth
= \tfrac{3\pi}{2}. \]
\end{example}

\begin{pitfallbox}
\begin{itemize}[leftmargin=*,nosep]
  \item $(r, \theta)$ is not unique: $(r, \theta) = (-r, \theta + \pi)$ and $(r, \theta + 2\pi k)$
        give the same point.
  \item Setting $r = f(\theta) = 0$ gives the values of $\theta$ where the curve passes through
        the pole --- useful for finding limits of integration.
  \item To find intersections of two polar curves, also check whether both pass through the pole;
        the system $r_1 = r_2$ may miss intersections at the origin.
\end{itemize}
\end{pitfallbox}
